\chapter{Tokens — the unit of work}
\label{ch:tokens}

\newthought{In Chapter~\ref{ch:what-lm}} we wrote ``each \(x_t\) is one
element of a vocabulary of size \(V\)'' and waved our hands. That hand wave
hides one of the most consequential design choices in the entire stack. This
chapter unpacks it.

\section{Why characters and why not}

The simplest possible vocabulary is the alphabet: 26 letters plus punctuation
and digits. Vocab size \(V\) becomes ${\sim}100$, the model never sees an
\term{out-of-vocabulary} word, and tokenisation is trivial.\sidenote{The
extreme of character-level modelling is byte-level: \(V = 256\), every
sequence of bytes is representable, no Unicode normalisation needed. Models
like ByT5 take this seriously.}

The price you pay is sequence length. ``the unfortunate consequences of
quadratic scaling'' is 7 words but about 50 characters. A character model
sees 7$\times$ as many tokens for the same input. Since the cost of attention
grows with the square of sequence length (the whole point of
Chapter~\ref{ch:why-quadratic}), that 7$\times$ becomes a 50$\times$ compute
penalty per layer.

\section{Why words and why not}

The opposite extreme is whole words: \(V \approx 10^5\) for English. Sequences
become short, but you suddenly have to handle words you have never seen
before — proper nouns, technical jargon, typos, every other language. Most
serious work has used \term{subwords} since around 2016.

\section{Subword tokenisation in one page}

The dominant tokenisers all do roughly the same thing: build a vocabulary by
merging frequent character pairs, then tokenise text greedily against that
vocabulary.

\paragraph{Byte-pair encoding (BPE).} Start with a vocabulary of all
characters. Find the most frequent adjacent pair in your corpus
(say, ``\code{t h}''), merge it into a single token (``\code{th}''), and add
that to the vocabulary. Repeat until you hit a target \(V\). This was
originally a compression algorithm; \citet{sennrich2016bpe} adapted it for
neural machine translation, and GPT-2 / GPT-3 popularised it for
LLMs.\sidenote{See \citep{sennrich2016bpe} for the original NLP application
and the OpenAI tokenizer notebooks for the engineering reality.}

\paragraph{WordPiece.} Used by BERT. Same merging idea, slightly different
scoring criterion (likelihood gain rather than raw frequency).\sidenote{Tied
to the \texttt{bert-base-uncased} vocabulary and the famous ``\code{\#\#}''
prefix that marks ``this is a continuation of the previous token''.}

\paragraph{SentencePiece (unigram).} Used by Llama, Mistral, and Gemma.
Treats tokenisation as a probabilistic model: pick the segmentation that
maximises a per-token probability. Operates directly on raw bytes and does
not require pre-tokenising into words first, which is a real win for
languages without whitespace boundaries.\sidenote{\citep{kudo2018sentencepiece}.
SentencePiece is the reason Llama-style models can tokenise Japanese without
a custom segmenter.}

\section{Why this affects everything downstream}

A subword tokeniser fixes a vocabulary \(V\) at training time. Three
consequences ripple through the rest of the system:

\begin{enumerate}
\item \textbf{The output projection costs \(\bigO(V \cdot d)\) per token.}
At each step, the model converts a hidden vector of dimension \(d\) into
\(V\) logits via a linear layer. If \(V\) is 128,000 and \(d\) is 4,096
that is ${\sim}500$ million multiply-adds per token, just for the last layer.

\item \textbf{Sequence length is in tokens, not characters.} A 1M-token
context window means roughly 750k English words, or 3,000 pages of
double-spaced text — assuming an English subword tokeniser. Code, Chinese,
or mixed content have different token-to-character ratios.\sidenote{Rule of
thumb for the OpenAI/Llama tokenisers: 1 token $\approx$ 4 English
characters $\approx$ 0.75 English words. Code is denser; one short
identifier can be one token.}

\item \textbf{The choice of tokeniser bakes in biases.} A tokeniser trained
mostly on English will fragment Hindi or Korean into many small tokens,
making non-English usage measurably more expensive both in compute and in
context window.
\end{enumerate}

\keyidea{Tokens are a cheap-looking choice with expensive consequences. The
``context length'' a vendor advertises is in their tokens, not yours, and
the shape of the cost curves in later chapters is in token-space.}

\section{Embeddings: turning tokens into vectors}

Once you have a token id (an integer in \([0, V)\)), the very first thing the
model does is look it up in a table:
\begin{equation}
  e_t \;=\; E[x_t] \in \R^{d},
\end{equation}
where \(E \in \R^{V \times d}\) is the \term{embedding matrix}, learned along
with everything else. Every transformer block sees only these dense vectors,
not the original token ids.

For a 7B parameter model with \(V = 32{,}000\) and \(d = 4{,}096\) the
embedding table alone is ${\sim}130$M parameters — roughly 2\% of the
model.\sidenote{Many models tie the embedding matrix to the output
projection (use the same \(E\) on the way in and the way out) which halves
this cost. Tied embeddings are standard in Llama; OpenAI's are not.}

\section{The same line, three tokenisers}

The following sentence (one paragraph of prose plus one line of Python) was
run through three public tokenisers. The exact byte sequence is:

\begin{quote}\small\ttfamily
The quick brown fox jumps over the lazy dog.\\
def fibonacci(n): return n if n < 2 else fibonacci(n-1) + fibonacci(n-2)
\end{quote}

\noindent
\textsc{Gpt2} (Byte-level BPE as shipped with GPT-2 — reproduced via
\texttt{tiktoken}) yields \(42\) tokens. Representative splits include
\texttt{fib}, \texttt{on}, \texttt{acci} rather than a single
\texttt{fibonacci} merge:

\begin{quote}\footnotesize\ttfamily\raggedright
The; quick; brown; fox; jumps; over; the; lazy; dog;.; \textbackslash n;
def; fib; on; acci; (; n; );:; return; n; if; n; <; 2; else; fib; on; acci; \ldots
\end{quote}

\noindent
\textsc{bert-base-uncased} (WordPiece) yields \(45\) tokens. The function name
splits as \texttt{fi} + \texttt{\#\#bon} + \texttt{\#\#ac} + \texttt{\#\#ci} — the
\texttt{\#\#} prefix marks a continuation piece:

\begin{quote}\footnotesize\ttfamily\raggedright
the; quick; brown; fox; jumps; over; the; lazy; dog; .; def; fi; \#\#bon; \#\#ac; \#\#ci; (; n; ); :; return; n; if; n; <; 2; else; fi; \ldots
\end{quote}

\noindent
\textsc{Llama-family SentencePiece} (TinyLlama checkpoint — same SPM pipeline as
Llama~2/3; Llama~3 weights are gated on HuggingFace so we use an open
checkpoint for reproducibility) yields \(45\) tokens. ``Fox'' can split as
\texttt{fo}+\texttt{x} depending on merge stats; \texttt{fibonacci} is still
broken into several pieces because the merge never entered the vocabulary as a
single token during training:

\begin{quote}\footnotesize\ttfamily\raggedright
The; quick; brown; fo; x; j; umps; over; the; lazy; dog; .; \textbackslash n;
def; fib; on; acci; (; n; );:; return; n; if; n; <; ; 2; else; fib; on; acci; \ldots
\end{quote}

\paragraph{Takeaway.} Code is not inherently ``denser'' than English — it depends
on whether the BPE merges exist for your identifiers. English prose shares many
merges across corpora; a rare function name is often glued together from
shorter atoms unless the tokenizer saw that exact string millions of times during
vocabulary construction.

\section{Back-of-the-envelope token economics}

Table~\ref{tab:token-econ} converts rough document sizes into token counts with
the OpenAI \texttt{cl100k\_base} convention (\(\approx\)~1.3~tokens per English
word for mixed prose). Dollar amounts use an \emph{illustrative} GPT-class input
price of USD~2.50 per million tokens — swap in your provider's current page when
you budget.\sidenote{Pricing moves quarterly; the point is order-of-magnitude,
not the third decimal.}

\begin{table}[htbp]
\centering
\small
\begin{tabular}{@{}lrrr@{}}
\toprule
Document & Approx.\ tokens & Minutes of GPU ``talk time''\,$^{*}$ & Illustr.\ USD @ \$2.50/M \\
\midrule
10k-word English essay & ${\sim}13{,}000$ & --- & ${\sim}\$0.03$ \\
10k-word Hindi essay & ${\sim}28{,}000$--$35{,}000$ & --- & ${\sim}\$0.07$--\$0.09 \\
1k-line Python file (${\sim}$40~char/line) & ${\sim}30{,}000$--$45{,}000$ & --- & ${\sim}\$0.08$--\$0.11 \\
100-page contract (${\sim}$450 words/page) & ${\sim}60{,}000$--$75{,}000$ & --- & ${\sim}\$0.15$--${\sim}\$0.19$ \\
\bottomrule
\end{tabular}\\[0.5em]
\footnotesize
$^{*}$GPU time is orthogonal to dollar pricing but scales with tokens \(\times\) depth.
\caption{Rule-of-thumb token loads for budgeting API spend (not for fitting KV caches).}
\label{tab:token-econ}
\end{table}

Non-English text routinely consumes \emph{more} tokens per unit of meaning when
the tokenizer was trained primarily on English — so equality-sensitive products
should quote prices \emph{per language} or risk silently taxing non-English
users.\sidenote{See also fragmentation studies on multilingual \texttt{cl100k}
and Llama tokenisers; the qualitative point is stable even when exact factors
shift.}

The next chapter is about how, before the transformer, people stitched these
embeddings together over time without an attention matrix at all.
